% !TeX root = ../main.tex

\chapter{Estado del arte}
\label{chap:estado_arte}

\section{Fundamentos de los modelos de lenguaje relevantes para este trabajo}
\label{sec:fundamentos_llm}

Los modelos de lenguaje utilizados en este trabajo pertenecen a la familia de modelos basados en la arquitectura \textit{Transformer}, introducida por Vaswani et al.~\cite{vaswani2017attention}. Esta arquitectura sustituye las dependencias recurrentes tradicionales por mecanismos de atención, permitiendo modelar relaciones entre los elementos de una secuencia con independencia de su distancia. Sobre esta base se han desarrollado los modelos de lenguaje autorregresivos modernos, que generan una respuesta estimando sucesivamente la distribución de probabilidad del siguiente token condicionada por el contexto precedente~\cite{brown2020language}.

Esta naturaleza probabilística resulta especialmente relevante para el diseño experimental del presente trabajo. Cuando la generación emplea estrategias de muestreo, una misma entrada no determina necesariamente una única continuación: la respuesta se obtiene seleccionando tokens de una distribución cuya forma depende tanto del modelo como de los parámetros de decodificación. Parámetros
como la temperatura y el \textit{top-p} controlan este proceso; este último, conocido también como \textit{nucleus sampling}, restringe el muestreo al conjunto de tokens que concentra una determinada masa de probabilidad~\cite{holtzman2019curious}. Estas variables deben mantenerse constantes cuando se comparan distintas condiciones experimentales. De ello se deriva una consecuencia directa para la medición: dos ejecuciones con una entrada idéntica pueden producir respuestas distintas, de modo que cualquier diferencia observada entre condiciones debe interpretarse frente a esa variabilidad basal y no frente a un cero teórico. Esta propiedad motiva las referencias de ruido empleadas en el diseño experimental, materializadas mediante la condición placebo y el agente ciego.

Otra propiedad central es la capacidad de los modelos de adaptar su comportamiento a la información proporcionada en el contexto. Brown et al.~\cite{brown2020language} mostraron que un modelo puede realizar nuevas tareas a partir de instrucciones y ejemplos incluidos directamente en la entrada, sin necesidad de actualizar sus parámetros. El condicionamiento contextual permite utilizar el \textit{prompt} no solo para describir la tarea, sino también para asignar funciones, proporcionar información específica o establecer restricciones sobre el comportamiento esperado. Además, se ha observado que los modelos pueden ser sensibles incluso a modificaciones del diseño del \textit{prompt} que preservan esencialmente su contenido semántico~\cite{sclar2024quantifying}. En este trabajo, esta capacidad de modificar el comportamiento mediante la información incluida en el \textit{prompt} sustenta tanto los distintos niveles de instrucción asignados a los agentes como las intervenciones de mitigación evaluadas posteriormente.

A esta característica se suma el desarrollo de modelos específicamente ajustados para seguir instrucciones. Procedimientos de ajuste supervisado y aprendizaje a partir de preferencias humanas han mostrado que el comportamiento de un modelo preentrenado puede orientarse hacia un mejor seguimiento de las intenciones expresadas por el usuario~\cite{ouyang2022training}. Este proceso de alineamiento depende de los ejemplos, preferencias y criterios utilizados durante el ajuste, por lo que no debe suponerse que todos los tipos de comportamiento o contenido hayan recibido una cobertura equivalente. La posibilidad de que dos atributos sensibles distintos hayan sido tratados con distinta intensidad durante el alineamiento se retoma en el Capítulo~\ref{chap:discusion}.

\section{Sesgo y \textit{fairness} en modelos de lenguaje}
\label{sec:sesgo_fairness}

El término \textit{sesgo} se utiliza en el ámbito del aprendizaje automático para describir fenómenos distintos y no necesariamente equivalentes. En el caso del procesamiento del lenguaje natural puede referirse, entre otros comportamientos, a la reproducción de estereotipos, a diferencias en el tratamiento de grupos sociales o a decisiones que favorecen sistemáticamente a unas categorías frente a otras. Esta diversidad hace que la presencia de sesgo no pueda determinarse de forma independiente de la tarea y del daño que se pretende estudiar. Blodgett et al.~\cite{blodgett2020language} señalan precisamente que una parte de la literatura sobre sesgo en procesamiento del lenguaje emplea el término sin explicitar suficientemente qué comportamiento se considera perjudicial, para quién y bajo qué criterio normativo.

Esta necesidad de concretar el objeto de análisis cobra especial importancia en los modelos de lenguaje de propósito general. Estos sistemas pueden manifestar asociaciones sociales tanto en el texto generado como en tareas en las que la salida constituye una valoración, recomendación o decisión. Gallegos et al.~\cite{gallegos2024bias} distinguen diferentes formas de operacionalizar el sesgo en LLM y muestran que su evaluación puede realizarse sobre representaciones internas, probabilidades del modelo o directamente sobre las respuestas generadas. En consecuencia, detectar una asociación estereotipada en el lenguaje y detectar un tratamiento diferencial en una tarea de decisión son problemas relacionados, pero no equivalentes.

En este trabajo interesa principalmente el segundo caso, el tratamiento diferencial en una tarea de asignación. La salida de cada agente no se evalúa por la presencia de una determinada expresión lingüística, sino por la asignación monetaria que realiza sobre un conjunto de empresas. El comportamiento de interés aparece cuando una modificación de información que debería permanecer ajena a los fundamentos financieros se asocia con un cambio en esa asignación. Esta definición traslada el estudio desde las asociaciones presentes en el lenguaje hacia sus posibles consecuencias en una tarea de asignación de recursos.

El concepto de \textit{fairness} introduce un criterio con el que comparar esas asignaciones. No existe una única definición de equidad aplicable a cualquier sistema, ya que distintas formulaciones imponen condiciones diferentes sobre qué propiedades de una asignación deberían mantenerse invariantes entre individuos o grupos~\cite{gallegos2024bias}. Para el diseño de este trabajo se toma como referencia la perspectiva contrafactual. Kusner et al.~\cite{kusner2017counterfactual} formalizan la
\textit{counterfactual fairness} mediante la idea de que una decisión relativa a un individuo debería permanecer invariante en escenarios contrafactuales en los que cambia un atributo sensible. Su formulación se apoya en un modelo causal estructural que permite definir dichos escenarios. El presente trabajo no pretende reproducir esa formalización causal completa, sino trasladar su intuición de invariancia a un entorno experimental controlado en el que los contrafactuales se construyen mediante intervención directa sobre los datos.

El sujeto de la decisión también difiere, los enfoques de \textit{fairness} contrafactual se plantean habitualmente sobre decisiones que afectan a un individuo perteneciente a un grupo protegido, mientras que aquí la unidad que recibe la asignación es una empresa. Los atributos estudiados son características asociadas a ella: el género de la persona que ocupa su dirección ejecutiva y la localización de su sede. El posible efecto adverso no se materializa, por tanto, como una decisión desfavorable sobre un individuo concreto, sino como un patrón de asignación de capital que favorezca o penalice sistemáticamente a empresas asociadas a determinadas características. Este tratamiento diferencial por asociación es el comportamiento que el diseño experimental pretende detectar.

Sobre esta base se construyen los pares contrafactuales utilizados posteriormente. Dos versiones de una misma empresa mantienen constantes sus fundamentos financieros y modifican únicamente el atributo objeto de estudio. Si las decisiones fueran invariantes respecto a esa modificación, ambas
versiones deberían producir comportamientos comparables, salvo por la variabilidad propia del proceso de generación descrita en la Sección~\ref{sec:fundamentos_llm}. La comparación contrafactual debe realizarse, por tanto, respecto a esa variabilidad basal y no exigir igualdad exacta entre dos ejecuciones concretas.

Esto delimita el alcance del término \textit{sesgo} en el resto de la memoria. El objetivo es estudiar una forma concreta de sensibilidad contrafactual en una tarea de inversión controlada: hasta qué punto decisiones con los mismos fundamentos financieros cambian cuando se modifica información sobre el género de la dirección ejecutiva o la localización corporativa. La forma de cuantificar esas diferencias y las limitaciones de los distintos paradigmas de medición se desarrollan en la sección siguiente.

\section{Medición del sesgo: paradigmas y limitaciones}
\label{sec:medicion_sesgo}

\subsection{Paradigmas de evaluación}
\label{subsec:paradigmas_medicion}

La forma de medir el sesgo depende del nivel del sistema en el que se sitúe la evaluación. Una primera familia de métodos opera sobre propiedades internas del modelo, como las asociaciones presentes en sus representaciones vectoriales. Estas métricas permiten detectar relaciones entre conceptos y grupos sociales sin necesidad de ejecutar una tarea aplicada, pero su capacidad para anticipar el comportamiento posterior del sistema es limitada. Goldfarb-Tarrant et al.~\cite{goldfarb2021intrinsic}, tras comparar métricas intrínsecas y extrínsecas en distintos modelos, tareas y lenguas, no encuentran una correlación fiable entre ambas categorías. Para estudiar las consecuencias del sesgo en una decisión concreta resulta, por tanto, más informativo observar directamente el comportamiento producido en esa tarea.

Una segunda aproximación utiliza bancos de pruebas construidos específicamente para provocar o detectar asociaciones estereotipadas. CrowS-Pairs~\cite{nangia2020crows} compara pares de oraciones que expresan en distinto grado un estereotipo; StereoSet~\cite{nadeem2021stereoset} evalúa preferencias entre asociaciones estereotipadas y alternativas; y BBQ~\cite{parrish2022bbq} plantea preguntas en contextos informativos y ambiguos para comprobar si el modelo recurre a estereotipos sociales. Estos instrumentos proporcionan escenarios controlados y una referencia explícita con la que clasificar las respuestas, por lo que resultan adecuados para estudiar la reproducción de estereotipos mediante salidas categóricas.

Su objetivo, sin embargo, es distinto del considerado en esta memoria. La decisión de inversión no contiene una respuesta categórica que pueda etiquetarse directamente como correcta, neutral o estereotipada. La salida es una cantidad monetaria continua cuyo posible tratamiento diferencial adquiere significado al compararla con la decisión obtenida para la misma empresa bajo otra condición. La ausencia de verbalización explícita del atributo tampoco garantiza que este no haya influido sobre la decisión.

Esta última distinción ha comenzado a estudiarse mediante evaluaciones centradas en las acciones de agentes basados en LLM. Li et al.~\cite{li2025actions} muestran que disparidades apenas visibles mediante preguntas explícitas pueden aparecer cuando el modelo debe adoptar decisiones con consecuencias sobre terceros. Su metodología traslada a agentes basados en LLM la lógica de los estudios de auditoría por correspondencia, en los que se comparan respuestas ante perfiles equivalentes que difieren en una característica demográfica. No obstante, su cuantificación se apoya en diferencias agregadas de tasas, mediante \textit{demographic parity difference}. Incluso en este paradigma orientado a decisiones, por tanto, el foco continúa situado principalmente en diferencias de nivel medio entre grupos y no en la estabilidad de las decisiones entre ejecuciones.

Las métricas tradicionales de \textit{fairness} de grupo constituyen otra familia de aproximaciones. Medidas como la paridad demográfica comparan tasas de decisión entre grupos, mientras que criterios como \textit{equalized odds} incorporan además la relación entre la decisión y un resultado observado. Estas métricas resultan apropiadas cuando existe una población dividida en grupos y una decisión que puede resumirse mediante tasas o categorías; varias de ellas requieren además una variable objetivo que actúe como referencia. Este supuesto difiere del escenario estudiado aquí, donde una misma empresa recibe una asignación monetaria continua y no existe una cantidad ``correcta'' de capital que permita determinar aciertos y errores.

El paradigma que mejor se adapta a este escenario es, en consecuencia, la comparación contrafactual pareada introducida en la Sección~\ref{sec:sesgo_fairness}. En lugar de comparar tasas agregadas entre dos poblaciones distintas, se contrastan las decisiones obtenidas para dos versiones de la misma unidad experimental, manteniendo constantes las variables que deben permanecer invariantes. La diferencia entre ambas asignaciones proporciona así una medida extrínseca y directamente vinculada a
la tarea.

\subsection{Variabilidad, arbitrariedad y referencia empírica}
\label{subsec:variabilidad_medicion}

Una comparación contrafactual pareada plantea todavía una cuestión adicional: qué propiedad de la distribución de las diferencias debe considerarse relevante. El enfoque más inmediato consiste en estudiar su media. Si una variante recibe sistemáticamente menos capital que su control, el \textit{gap} medio será negativo y permitirá identificar una dirección común del tratamiento diferencial. Sin embargo, dos condiciones pueden presentar medias semejantes y diferir sustancialmente en la estabilidad de las decisiones que producen.

Este problema tiene un antecedente en la literatura sobre arbitrariedad en la predicción. Cooper et al.~\cite{cooper2024arbitrariness} estudian cómo modelos igualmente admisibles obtenidos mediante un mismo procedimiento pueden producir decisiones diferentes para los mismos individuos. En particular,
denominan \textit{arbitrariedad sistemática} a las situaciones en las que esta inconsistencia no se distribuye de manera uniforme, sino que afecta en mayor medida a determinados grupos. Su conclusión principal va, de hecho, más allá: al contabilizar la arbitrariedad, la mayoría de los bancos de pruebas habituales resultan próximos a la equidad según las métricas empleadas, lo que sugiere que parte de la desigualdad documentada puede reflejar variabilidad no modelada. Al mismo tiempo, muestran que satisfacer una métrica agregada de \textit{fairness} no implica necesariamente que las decisiones sean igualmente consistentes para todos los casos. Ambas observaciones apuntan a la misma limitación: la equidad agregada y la estabilidad de las decisiones son dimensiones distintas, y evaluar únicamente la primera puede ofrecer una caracterización incompleta del sistema.

El origen de la arbitrariedad estudiada por Cooper et al. y el considerado en esta memoria es distinto. En su trabajo, la variabilidad aparece entre modelos alternativos que podrían haberse obtenido y desplegado a partir de un mismo procedimiento de aprendizaje. Aquí se estudian ejecuciones repetidas de un mismo modelo y una misma configuración. La arbitrariedad puede aparecer, por tanto, dentro de un sistema ya fijado, sin depender de qué modelo concreto haya seleccionado previamente el desarrollador. La conexión entre ambos problemas reside en que una medida agregada favorable puede ocultar que unas condiciones produzcan decisiones considerablemente menos reproducibles que otras.

En los modelos de lenguaje esta cuestión adquiere una relevancia adicional debido al carácter estocástico de la generación. Song et al.~\cite{song2025good} señalan que numerosas evaluaciones se basan en una única generación por ejemplo y que la variabilidad entre generaciones no suele reflejarse mediante estadísticas como la desviación típica. En esas condiciones, dos sistemas con un rendimiento medio semejante pueden presentar grados de estabilidad muy distintos sin que esa diferencia quede reflejada en la métrica principal de evaluación.

Cabe distinguir, por tanto, dos formas de tratamiento diferencial. La primera es un \textit{desplazamiento}: la presencia del atributo se asocia de forma consistente con una decisión mayor o menor que la obtenida en control. La segunda es un aumento de arbitrariedad: el promedio puede permanecer próximo al de control mientras las decisiones se vuelven menos estables cuando el
atributo está presente. Por analogía con la terminología de Cooper et al., un aumento sistemático de esta segunda componente bajo determinadas condiciones puede caracterizarse como una forma de arbitrariedad sistemática, aunque aquí la comparación se realiza entre condiciones contrafactuales y no entre grupos de individuos.

Esta distinción tiene consecuencias prácticas para la evaluación de equidad. Un procedimiento de auditoría centrado exclusivamente en diferencias de medias o tasas podría considerar satisfactorio un sistema en el que ambos grupos reciben, en promedio, un tratamiento semejante, aun cuando uno de ellos esté sujeto a decisiones mucho menos reproducibles. La relevancia de esta dimensión para los marcos de evaluación y supervisión de sistemas de IA se retoma en el Capítulo~\ref{chap:marco_regulador}.

La existencia de variabilidad basal introduce además un problema para definir la referencia de comparación. Como se explicó en la Sección~\ref{sec:fundamentos_llm}, dos ejecuciones idénticas no tienen por qué producir exactamente la misma asignación. Por ello, un \textit{gap} distinto de cero en una pareja concreta no basta para atribuir la diferencia al atributo manipulado. La magnitud observada debe compararse con la variabilidad que el propio sistema produce cuando no existe una diferencia informativa entre las condiciones.

Esta cuestión resulta particularmente importante al estudiar la dispersión. Una varianza elevada de los \textit{gaps} no constituye por sí sola evidencia de sensibilidad al atributo, ya que parte de ella puede proceder del propio proceso de generación. Es necesario estimar la variabilidad basal del sistema y determinar si la condición estudiada presenta un exceso respecto a esa referencia. El diseño desarrollado posteriormente incorpora dos controles complementarios con esta finalidad: una condición placebo formada por pares idénticos como referencia externa de ruido y un agente ciego al atributo como referencia interna sobre las mismas parejas experimentales.

La evaluación simultánea de múltiples modelos, niveles de instrucción y atributos introduce además un problema de multiplicidad. Cuando se realizan numerosos contrastes sobre una misma pregunta experimental, interpretar cada valor $p$ de forma aislada incrementa la probabilidad de identificar
hallazgos espurios. El procedimiento de Benjamini y Hochberg~\cite{benjamini1995controlling} permite controlar la tasa esperada de falsos descubrimientos entre las hipótesis rechazadas. Por este motivo, los análisis posteriores definen explícitamente las familias de contrastes y aplican esta corrección antes de considerar un resultado como evidencia inferencial.

En conjunto, el marco de medición adoptado combina una comparación contrafactual sobre el resultado efectivo de la decisión con dos dimensiones complementarias: el desplazamiento medio y la dispersión de los \textit{gaps}. Ambas se interpretan respecto a referencias empíricas de variabilidad basal y dentro de familias de contrastes previamente definidas. Las métricas y procedimientos estadísticos concretos utilizados para operacionalizar estas ideas se detallan en el Capítulo~\ref{chap:metodologia}.

\section{Sistemas multiagente y dinámica deliberativa}
\label{sec:mas_deliberacion}

La utilización de modelos de lenguaje como agentes permite construir sistemas en los que varias instancias reciben una tarea, mantienen estados o contextos propios y se comunican para alcanzar un resultado conjunto. A diferencia de una estrategia de muestreo independiente, en la que varias respuestas se generan sin interacción y se agregan posteriormente, un sistema multiagente introduce dependencias entre las decisiones: la respuesta producida por un agente puede convertirse en parte del contexto que condiciona las respuestas posteriores de los demás. El comportamiento del sistema depende, por tanto, no solo de los modelos que lo componen, sino también de cómo se definen sus
roles y de la estructura de comunicación empleada.

Una de las formas habituales de especializar estos sistemas consiste en asignar funciones mediante el propio \textit{prompt}. El marco CAMEL~\cite{li2023camel}, por ejemplo, utiliza \textit{role-playing} para condicionar a distintos agentes hacia funciones complementarias y facilitar su cooperación. Esta especialización no constituye únicamente una descripción superficial del agente: las instrucciones que definen su función modifican el contexto desde el que interpreta la tarea y genera sus respuestas. En esta línea, ChatEval muestra que la diversidad de roles puede
afectar al resultado de la deliberación y que utilizar descripciones de rol idénticas para todos los agentes puede degradar el rendimiento del sistema~\cite{chan2024chateval}. Estos antecedentes sustentan el uso, en este trabajo, de distintos grados de especificidad en los roles asignados a los
miembros del comité.

La interacción entre agentes se ha estudiado especialmente mediante protocolos de debate. Du et al.~\cite{du2023improving} proponen generar primero respuestas individuales y permitir después que varias instancias examinen las respuestas y razonamientos de las demás durante sucesivas rondas. En sus experimentos, este procedimiento mejora el razonamiento y la factualidad respecto al uso de una única instancia. Liang et al.~\cite{liang2024encouraging} muestran igualmente que el intercambio entre agentes puede favorecer la aparición de alternativas que un modelo aislado no genera mediante autorreflexión. La deliberación ofrece así un mecanismo para incorporar información procedente de distintas respuestas iniciales y revisar decisiones previas.

Esta mejora no constituye, sin embargo, una propiedad garantizada de cualquier sistema deliberativo. El resultado depende de factores como las opiniones iniciales, la composición del grupo y la forma en que los agentes actualizan sus respuestas. Trabajos recientes sobre debate multiagente muestran que los protocolos simples no superan necesariamente a mecanismos de agregación más sencillos y señalan la diversidad de las respuestas iniciales como uno de los factores que condicionan el beneficio de la interacción \cite{zhu2026demystifying}. Por tanto, incorporar más agentes o más rondas de comunicación no puede suponerse a priori como una mejora del comportamiento del sistema.

La composición del comité constituye en este contexto una variable propia del sistema. Un comité homogéneo replica un mismo modelo en los distintos agentes, aunque estos puedan recibir contextos o roles diferentes, mientras que un comité heterogéneo combina modelos distintos dentro del mismo proceso deliberativo. Estas dos configuraciones pueden generar grados diferentes de diversidad inicial y de dependencia entre las respuestas. La literatura sobre debate proporciona evidencia de que dicha diversidad puede modificar el resultado de la interacción~\cite{chan2024chateval,zhu2026demystifying}, aunque su efecto depende de la tarea y del protocolo empleado. Esta distinción motiva que la composición se trate en este trabajo como una dimensión experimental y que los comités homogéneos y heterogéneos se evalúen bajo un mismo procedimiento de deliberación.

La comunicación hace además que el sistema deba entenderse como un proceso dinámico y no únicamente por su salida final. Chuang et al.~\cite{chuang2024simulating} muestran, en redes de agentes basados en LLM, que el intercambio iterativo puede producir dinámicas diferentes de convergencia o fragmentación dependiendo de cómo estén condicionados los agentes. Cada ronda modifica el contexto disponible y puede alterar tanto la posición individual como el grado de acuerdo entre los participantes. La interacción puede reducir diferencias iniciales, pero también generar dependencias entre respuestas y hacer que una posición o error previamente localizado influya sobre otros miembros del sistema.

Estas dependencias implican que la información disponible para cada agente deja de estar determinada únicamente por su entrada inicial. Un miembro del comité que no reciba directamente un dato concreto puede verse influido posteriormente por él si otros agentes sí lo han observado y lo incorporan a sus decisiones o justificaciones. La estructura de comunicación determina, por tanto, no solo qué información recibe inicialmente cada agente, sino también qué información puede alcanzarlo de forma indirecta durante la deliberación. Wu et al.~\cite{wu2025can} muestran, en este sentido, que la presión mayoritaria puede dificultar la corrección independiente y que los consensos incorrectos no siempre son descartados durante la interacción. Zhu et al.~\cite{zhu2026demystifying} señalan igualmente que las contribuciones de baja calidad pueden propagarse entre rondas en lugar de ser eliminadas por el grupo.

Esta posibilidad tiene una consecuencia directa para el diseño experimental del presente trabajo. Mantener uno de los agentes ciego al atributo estudiado permite observar si su comportamiento permanece independiente mientras no existe comunicación y cómo cambia una vez comienza a recibir las respuestas de los agentes que sí disponen de esa información. El agente ciego funciona así como una sonda de la posible transmisión indirecta de información durante la deliberación.

Esta perspectiva resulta también relevante para la estabilidad de las decisiones. Si las respuestas iniciales contienen variación parcialmente independiente, el intercambio de información puede favorecer su convergencia; si, por el contrario, los agentes adoptan o refuerzan las mismas señales, la comunicación puede aumentar la dependencia entre sus decisiones. Por ello, la evolución de un sistema multiagente no debe caracterizarse únicamente mediante el cambio de su respuesta media o mediante la decisión alcanzada al final del debate, sino también mediante la evolución de la variabilidad y del acoplamiento entre sus miembros.

En consecuencia, el protocolo empleado en este trabajo conserva una decisión inicial obtenida antes de cualquier comunicación y registra posteriormente la respuesta de cada agente durante cuatro rondas de interacción. Se utiliza un número fijo de turnos y no un criterio de parada basado en consenso, de modo que todas las condiciones experimentales compartan la misma longitud deliberativa y puedan compararse en cada momento del proceso. Esta estructura permite distinguir el comportamiento previo a la interacción de los cambios que aparecen una vez los agentes comienzan a observarse entre sí. La sección siguiente examina específicamente cómo estas dinámicas pueden afectar al sesgo en sistemas multiagente.

\section{Sesgo emergente en sistemas multiagente}
\label{sec:sesgo_mas}

La interacción entre agentes introduce una dimensión adicional en el estudio del sesgo. Aunque cada miembro de un sistema multiagente esté basado en un modelo cuyo comportamiento pueda evaluarse individualmente, la comunicación entre miembros y los mecanismos mediante los que se agregan sus respuestas generan un proceso colectivo cuyo resultado no tiene por qué coincidir con el de ninguno de sus componentes. En este contexto, el término \textit{sesgo emergente} puede adoptar dos significados relacionados. Por un lado, puede describir un comportamiento del sistema que no resulta directamente predecible a partir del sesgo de los agentes que lo componen. Por otro, puede referirse a la aparición temporal de un comportamiento sesgado que no estaba presente en fases anteriores de la interacción.

Los primeros trabajos sobre interacción social entre agentes basados en LLM ya muestran que la deliberación no elimina necesariamente las asociaciones aprendidas por los modelos. Taubenfeld et al.~\cite{taubenfeld2024systematic}, al estudiar simulaciones de debates políticos, observan que los agentes tienden a reproducir sesgos sociales del modelo subyacente incluso cuando se les asignan perspectivas específicas desde las que deben argumentar. El rol proporcionado mediante el \textit{prompt} no garantiza, por tanto, que las asociaciones del modelo dejen de influir en su comportamiento.

Borah y Mihalcea~\cite{borah2024towards} estudian específicamente el sesgo implícito de género durante interacciones multiagente. Sus experimentos muestran asociaciones sesgadas en una proporción elevada de los escenarios evaluados y, además, que estas asociaciones pueden intensificarse después de la interacción entre agentes. El resultado es relevante porque sitúa el problema no únicamente en el sesgo contenido inicialmente por cada modelo, sino también en la dinámica mediante la que varias respuestas se condicionan mutuamente. La interacción puede actuar así como un mecanismo que preserve o refuerce diferencias ya presentes en los agentes.

Un análisis más explícito de esta dimensión temporal es el desarrollado por Nguyen et al.~\cite{nguyen2025social}. Los autores descomponen una interacción multiagente en una fase de \textit{Genesis}, en la que cada agente genera una respuesta de forma independiente antes de recibir información de sus compañeros, y una fase posterior de comunicación iterativa. Sobre esta estructura distinguen tres fenómenos: la \textit{emergencia}, cuando aparece por primera vez una respuesta sesgada; la \textit{propagación}, cuando dicho comportamiento alcanza a otros agentes mediante la comunicación; y la \textit{amplificación}, cuando su presencia aumenta a lo largo de la interacción.

Sus experimentos sobre tres bancos de pruebas de estereotipos muestran que estas dinámicas pueden alcanzar magnitudes importantes. Dependiendo de la configuración, la comunicación llega a provocar hasta un 70\,\% de nueva emergencia de sesgo, permite extender respuestas sesgadas a más del 80\,\% de los agentes y produce amplificaciones superiores a tres veces el nivel inicial \cite{nguyen2025social}. La emergencia, sin embargo, se concentra en las dos primeras etapas: alrededor de la mitad de los casos aparece ya durante \textit{Genesis}, antes de cualquier intercambio, y casi todo el resto se produce en el primer turno de comunicación. Las rondas posteriores apenas introducen casos nuevos y actúan sobre todo estabilizando comportamientos ya establecidos. Una parte sustancial del fenómeno precede así a cualquier influencia social, lo que proporciona un antecedente directo para analizar el estado inicial y la deliberación de forma diferenciada.

En la mayoría de las condiciones de Nguyen et al., los agentes reciben roles asociados a grupos sociales relacionados con el estereotipo evaluado, de modo que la identidad de grupo forma parte de la propia definición del agente. Su diseño incluye, sin embargo, una condición neutral o \textit{Ungrouped}, sin asignación de grupo social. Incluso en este régimen, los sistemas multiagente muestran en numerosas configuraciones una robustez inferior a la de los agentes individuales, lo que indica que la pérdida de robustez no depende exclusivamente de introducir el atributo social mediante el rol \cite{nguyen2025social}.

El trabajo de Madigan et al.~\cite{madigan2025emergent} aborda el problema desde una perspectiva diferente y especialmente próxima al dominio financiero. Los autores construyen sistemas de tres agentes basados en distintos LLM y los evalúan en tareas de clasificación tabular relacionadas con ingreso y riesgo crediticio. El comportamiento del sistema multiagente se compara con el de sus modelos constituyentes funcionando individualmente, lo que permite estudiar si las propiedades de equidad del conjunto pueden inferirse a partir de las de sus componentes.

Los resultados muestran que esta relación no es estable. En términos agregados, la mediana del cambio resulta ligeramente favorable en la mayoría de métricas, con la diferencia de precisión como excepción, de modo que los sistemas multiagente tienden a reducir modestamente el sesgo. Sin embargo, las distribuciones presentan colas positivas pronunciadas y existen configuraciones en las que el sistema amplifica el sesgo hasta aproximadamente diez veces respecto a alguno de sus componentes \cite{madigan2025emergent}. Una tendencia promedio favorable puede coexistir, por tanto, con configuraciones concretas que producen incrementos muy elevados. Este resultado refuerza la necesidad de evaluar el sistema completo en lugar de inferir sus propiedades a partir del comportamiento aislado de los modelos que lo forman.

El diseño del presente trabajo se separa de estos antecedentes en varios aspectos. El primero afecta a la naturaleza del atributo evaluado. Aquí los roles de los agentes son profesionales y no guardan relación con el género de la persona que ocupa la dirección ejecutiva ni con la localización de la sede: el atributo pertenece a la empresa sobre la que se toma la decisión y no a la identidad de los miembros del comité.

El segundo se refiere a la variable manipulada. Madigan et al. comparan sistemas multiagente heterogéneos con sus modelos constituyentes funcionando como agentes individuales. Aquí se mantiene constante la existencia del comité y el protocolo de interacción, y se comparan composiciones homogéneas y heterogéneas. De este modo, la variable modificada es la diversidad de modelos dentro del sistema sin confundirla con la presencia o ausencia de deliberación.

También difieren los mecanismos mediante los que se operacionaliza el sesgo. Nguyen et al. utilizan tareas de respuesta múltiple en las que existen alternativas estereotipadas y una respuesta neutral identificable, mientras que Madigan et al. estudian clasificación binaria con etiquetas observadas y métricas de \textit{fairness} entre grupos. En ambos casos el análisis se apoya principalmente en resultados categóricos. El presente trabajo considera, en cambio, una asignación monetaria continua sin una cantidad correcta de capital y utiliza pares contrafactuales para aislar el efecto del atributo estudiado.

Finalmente, los trabajos anteriores caracterizan el sesgo principalmente a través de tasas o diferencias agregadas entre grupos. No consideran la dispersión entre ejecuciones repetidas de una misma configuración como desenlace de equidad. Como se desarrolló en la Sección~\ref{subsec:variabilidad_medicion}, esta dimensión permite detectar situaciones en las que el tratamiento medio permanece estable pero la presencia del atributo se asocia con una mayor dispersión de las decisiones: unas ejecuciones pueden favorecer notablemente a la variante mientras otras la penalizan en magnitud similar, haciendo que ambos efectos se compensen al promediarlos.

El diseño experimental de esta memoria atiende por ello a las dos acepciones planteadas al inicio. Conservar separadamente la respuesta de génesis y las rondas posteriores permite situar en el tiempo la aparición del efecto y observar su evolución durante la deliberación. Comparar composiciones bajo un mismo protocolo, con las decisiones registradas a nivel de agente, permite examinar hasta qué punto el comportamiento del conjunto puede anticiparse desde el de sus miembros. La presencia de un agente sin acceso directo al atributo añade además una vía para estudiar si una diferencia inicialmente localizada en los agentes visibles alcanza después a un miembro que solo recibe esa información de forma indirecta.


\section{Aplicaciones financieras de modelos de lenguaje}
\label{sec:aplicaciones_financieras}

\section{Mitigación de sesgo}
\label{sec:estado_arte_mitigacion}

\section{Limitaciones de los trabajos existentes y posicionamiento}
\label{sec:posicionamiento}